\documentclass[10pt]{article}
\usepackage[preprint]{tmlr}
\usepackage{amsmath,amssymb,booktabs,graphicx,hyperref,url}
\hypersetup{hidelinks}\raggedbottom
\def\month{09}\def\year{2026}
\newcommand{\evidenceroot}{..}
\newcommand{\rv}[1]{\csname rv@#1\endcsname}
% Generated; do not edit.
\expandafter\def\csname rv@c32.behavior_only.day_ap\endcsname{0.025238}
\expandafter\def\csname rv@c32.behavior_only.day_roc\endcsname{0.820}
\expandafter\def\csname rv@c32.behavior_only.user_roc\endcsname{0.825}
\expandafter\def\csname rv@c32.fold.prevalence\endcsname{0.000714}
\expandafter\def\csname rv@c32.full.day_ap\endcsname{0.004859}
\expandafter\def\csname rv@c32.full.day_roc\endcsname{0.540}
\expandafter\def\csname rv@c32.full.user_roc\endcsname{0.589}
\expandafter\def\csname rv@c32.gain\endcsname{0.280}
\expandafter\def\csname rv@c32.gain.ci\endcsname{[0.204, 0.357]}
\expandafter\def\csname rv@c32.pooled.prevalence\endcsname{0.032071}
\expandafter\def\csname rv@c32.profile_only.day_ap\endcsname{0.001353}
\expandafter\def\csname rv@c32.profile_only.day_roc\endcsname{0.482}
\expandafter\def\csname rv@c32.profile_only.user_roc\endcsname{0.476}
\expandafter\def\csname rv@claim.C01\endcsname{verified with scope limits}
\expandafter\def\csname rv@claim.C02\endcsname{verified with scope limits}
\expandafter\def\csname rv@claim.C03\endcsname{provenance qualified}
\expandafter\def\csname rv@claim.C04\endcsname{provenance qualified}
\expandafter\def\csname rv@claim.C05\endcsname{verified with scope limits}
\expandafter\def\csname rv@claim.C06\endcsname{verified with scope limits}
\expandafter\def\csname rv@claim.C07\endcsname{verified with scope limits}
\expandafter\def\csname rv@claim.C08\endcsname{verified with scope limits}
\expandafter\def\csname rv@claim.C09\endcsname{verified with scope limits}
\expandafter\def\csname rv@claim.C10\endcsname{fixed input only}
\expandafter\def\csname rv@claim.C11\endcsname{pending dependency repair}
\expandafter\def\csname rv@claim.C12\endcsname{pending dependency repair}
\expandafter\def\csname rv@claim.C13\endcsname{pending dependency repair}
\expandafter\def\csname rv@claim.C14\endcsname{pending dependency repair}
\expandafter\def\csname rv@claim.C15\endcsname{pending dependency repair}
\expandafter\def\csname rv@claim.C16\endcsname{pending dependency repair}
\expandafter\def\csname rv@claim.C17\endcsname{pending dependency repair}
\expandafter\def\csname rv@claim.C18\endcsname{pending dependency repair}
\expandafter\def\csname rv@claim.C19\endcsname{conditional only}
\expandafter\def\csname rv@claim.C20\endcsname{pending dependency repair}
\expandafter\def\csname rv@claim.C21\endcsname{pending dependency repair}
\expandafter\def\csname rv@claim.C22\endcsname{pending dependency repair}
\expandafter\def\csname rv@claim.C23\endcsname{pending dependency repair}
\expandafter\def\csname rv@claim.C24\endcsname{pending dependency repair}
\expandafter\def\csname rv@claim.C25\endcsname{pending dependency repair}
\expandafter\def\csname rv@claim.C26\endcsname{pending dependency repair}
\expandafter\def\csname rv@claim.C27\endcsname{pending dependency repair}
\expandafter\def\csname rv@claim.C28\endcsname{verified with scope limits}
\expandafter\def\csname rv@claim.C29\endcsname{pending dependency repair}
\expandafter\def\csname rv@claim.C30\endcsname{pending dependency repair}
\expandafter\def\csname rv@claim.C31\endcsname{verified with scope limits}
\expandafter\def\csname rv@claim.C32\endcsname{verified with scope limits}
\expandafter\def\csname rv@claim.C33\endcsname{verified with scope limits}
\expandafter\def\csname rv@claim.C34\endcsname{verified with scope limits}
\expandafter\def\csname rv@claim.C35\endcsname{diagnostic only}
\expandafter\def\csname rv@claim.C36\endcsname{fixed input only}
\expandafter\def\csname rv@claim.C37\endcsname{pending dependency repair}
\expandafter\def\csname rv@claim.C38\endcsname{conditional only}
\expandafter\def\csname rv@claim.C39\endcsname{fixed input only}
\expandafter\def\csname rv@claim.C40\endcsname{outside label repair}
\expandafter\def\csname rv@claim.C41\endcsname{outside label repair}
\expandafter\def\csname rv@claim.C42\endcsname{outside label repair}
\expandafter\def\csname rv@claim.C43\endcsname{conditional only}
\expandafter\def\csname rv@claim.C44\endcsname{synchronized recovery editions}
\expandafter\def\csname rv@claim.C45\endcsname{pending dependency repair}
\expandafter\def\csname rv@claim.C46\endcsname{pending dependency repair}
\expandafter\def\csname rv@factorial.3b.full.day_roc\endcsname{0.683}
\expandafter\def\csname rv@factorial.3b.full.user_roc\endcsname{0.778}
\expandafter\def\csname rv@factorial.3b.no_profile.day_roc\endcsname{0.894}
\expandafter\def\csname rv@factorial.3b.no_profile.user_roc\endcsname{0.950}
\expandafter\def\csname rv@factorial.3b.no_psy.day_roc\endcsname{0.799}
\expandafter\def\csname rv@factorial.3b.no_psy.user_roc\endcsname{0.901}
\expandafter\def\csname rv@factorial.3b.shuffle_profile.day_roc\endcsname{0.432}
\expandafter\def\csname rv@factorial.3b.shuffle_profile.user_roc\endcsname{0.550}
\expandafter\def\csname rv@factorial.8b.full.day_roc\endcsname{0.412}
\expandafter\def\csname rv@factorial.8b.full.user_roc\endcsname{0.472}
\expandafter\def\csname rv@factorial.8b.no_profile.day_roc\endcsname{0.899}
\expandafter\def\csname rv@factorial.8b.no_profile.user_roc\endcsname{0.902}
\expandafter\def\csname rv@factorial.8b.no_psy.day_roc\endcsname{0.791}
\expandafter\def\csname rv@factorial.8b.no_psy.user_roc\endcsname{0.901}
\expandafter\def\csname rv@factorial.8b.shuffle_profile.day_roc\endcsname{0.394}
\expandafter\def\csname rv@factorial.8b.shuffle_profile.user_roc\endcsname{0.438}
\expandafter\def\csname rv@r42.all_users.day_ap\endcsname{0.011939}
\expandafter\def\csname rv@r42.all_users.day_roc\endcsname{0.960}
\expandafter\def\csname rv@r42.all_users.user_roc\endcsname{0.965}
\expandafter\def\csname rv@r42.changed_rows\endcsname{407}
\expandafter\def\csname rv@r42.my_positive_days\endcsname{1,362}
\expandafter\def\csname rv@r42.my_positive_users\endcsname{70}
\expandafter\def\csname rv@r42.rows\endcsname{330,295}
\expandafter\def\csname rv@r42.unseen_users.day_ap\endcsname{0.004978}
\expandafter\def\csname rv@r42.unseen_users.day_roc\endcsname{0.538}
\expandafter\def\csname rv@r42.unseen_users.user_roc\endcsname{0.589}
\expandafter\def\csname rv@r62.all_users.day_ap\endcsname{0.000585}
\expandafter\def\csname rv@r62.all_users.day_roc\endcsname{0.956}
\expandafter\def\csname rv@r62.all_users.user_roc\endcsname{0.962}
\expandafter\def\csname rv@r62.changed_rows\endcsname{23}
\expandafter\def\csname rv@r62.my_positive_days\endcsname{73}
\expandafter\def\csname rv@r62.my_positive_users\endcsname{5}
\expandafter\def\csname rv@r62.rows\endcsname{1,393,297}
\expandafter\def\csname rv@r62.unseen_users.day_ap\endcsname{0.000123}
\expandafter\def\csname rv@r62.unseen_users.day_roc\endcsname{0.583}
\expandafter\def\csname rv@r62.unseen_users.user_roc\endcsname{0.649}

\title{What Does a Benign-Trained Anomaly Detector Actually Detect?\\Complete Reading Edition: Corrected Labels and Current Evidence}
\author{\name Suramya Angdembay \email suramya.angdembay@usm.edu\\\addr University of Southern Mississippi
\AND\name Haiyan Tian \email haiyan.tian@usm.edu\\\addr University of Southern Mississippi}
\begin{document}
\maketitle
\begin{abstract}
A detector can rank unusual days without using the information we intended.
We study language models adapted to employee logs and separate three questions:
what enters the anomaly score, what information the model uses to predict
behavior, and what changes when we edit its internal features. Corrected CERT
labels preserve the main adapters but change evaluation and some experiment
choices. The strongest complete result is on r4.2: scoring only behavior targets
raises day ROC from \rv{c32.full.day_roc} to \rv{c32.behavior_only.day_roc} and average precision from \rv{c32.full.day_ap} to \rv{c32.behavior_only.day_ap}.
The model still sees the user's profile. Our internal audit remains part of the
method, but the corrected pilot does not yet establish an attack-specific
mechanism. This edition follows the argument in order and places implementation
details, provenance, and unresolved evidence in the appendix.
\end{abstract}
\section{The problem}
Organizations record logons, file activity, web activity, and email counts.
An insider threat involves misuse of legitimate access. Because such incidents
are rare, an anomaly detector learns regular activity and assigns larger scores
to days that look unusual. A language model can do this by treating the daily
measurements as structured text and asking how predictable its tokens are.

Our question is what that unpredictability represents. It might concern the
actions of a user, the user's profile, or whether the model has seen that user
during training. A good ranking score alone cannot distinguish these
possibilities. We first define the data and evaluation, then open the model,
and finally separate supported findings from the remaining causal questions.
The technical appendix carries the additional detail of the main manuscript.
\section{Data, labels, and populations}
\label{sec:data}
% claims: C01 C02 C03 C04
We study serialized employee activity from CERT r4.2 and r6.2
\citep{certdataset,glasser2013bridging,le2020analyzing}. CERT is a synthetic
insider-threat benchmark. A record is one user's observed day, assembled from
extracted sessions. The \texttt{DAY} line contains week and organizational
fields, \texttt{PSY} contains simulated personality scores, and \texttt{SES}
lines contain session activity counts. \texttt{SESCOUNT} records the number of
sessions. These are structured measurements written as text.

Our primary target is membership in the official master's incident window:
a day is positive if its recorded calendar interval intersects the user's
half-open incident interval $[\mathrm{start},\mathrm{end})$. Evaluation is
restricted to days with extracted sessions. This labels an \emph{incident-window
day}, not every action on that day. An insider's other days are called
\emph{window-negative}; they are not verified benign activity.

An audit found an incorrect calendar join in the previous evaluation labels.
The corrected day-index origin is December 28, 2009. Independent enumeration
reproduced every corrected label, and the original texts, ordering, and splits
were preserved. The resulting observed populations are:
\begin{center}\input{\evidenceroot/tables/recovery_populations.tex}\end{center}
The r4.2 changes comprise 177 positive-to-negative and 230 negative-to-positive
labels; r6.2 has 10 and 13, respectively. The master windows include 530 r4.2
and 26 r6.2 additional days with no extracted session, which are outside this
observation domain. Named malicious-event days define a different target; they
are retained as sensitivity labels and are not silently substituted for the
master-window target. The r6.2 master/detail discrepancies and one repaired
malformed master date are documented in the label audit.

\paragraph{Training exposure.}
The main LM adapters remain usable: their actual training and validation
user lists exclude all 70 r4.2 and all 5 r6.2 official insiders. This was
checked independently of the faulty day-label join. The main training/validation
counts are 851/79 non-insider users for r4.2 and 3,590/405 for r6.2. The
unseen-user evaluation pools contain all official insiders plus those validation
users: 42,468 days from 149 users and 142,072 days from 410 users.
``Unseen'' means absent from adaptation; these repeatedly analyzed validation
users are not a newly untouched test set. The factorial and smaller-model
adapters are separate training runs, with their own provenance qualifications.

\paragraph{Why labels affect more than a metric.}
Correcting labels can change a feature's rank, the sample used to fit a sparse
dictionary, or which user is eligible as a normal donor. In those cases the
historical intervention remains a measurement of its original inputs, but it
is not the experiment the corrected procedure would have selected. Model
training is reusable only when its actual inputs and fitting membership are
unchanged. We apply this distinction throughout.

\section{What the score measures}
\label{sec:evaluation}
% claims: C05 C06 C07 C08 C09 C31 C32 C33 C34
For tokens $x_1,\ldots,x_n$, the adapted LM assigns the anomaly score
\begin{equation}
 s(x)=-\frac{1}{n-1}\sum_{t=2}^{n}\log p_{\mathrm{ada}}(x_t\mid x_{<t}).
\end{equation}
The first token is context, padding is excluded, and losses are assigned to the
\emph{predicted target token}. A larger score means the observed day was less
predictable under the model of the adaptation corpus.

We separate two questions. In the general-population comparison, the negative
user pool includes users seen during adaptation. In the unseen-user comparison,
negative users are restricted to validation users absent from adaptation. Each
fold contains one official insider and a fixed sampled negative-user cohort;
metrics are averaged equally over insiders. The general comparison uses 800
non-insider users per fold; the unseen comparison uses all 79 or 405 validation
users. An insider's window-negative days are negatives for day metrics. User
scores are the maximum over each user's full scored history. The model stays
fixed across these evaluation folds.

ROC-AUC measures the probability that a positive example outranks a negative
one, with half credit for ties. Average precision (AP) summarizes precision
across recall levels. It is not precision at a chosen alert threshold. A fold
average and a pooled metric have different weighting and prevalence: neither
can be substituted for the other. Unless explicitly labeled pooled, the tables
below report equal-insider fold means. Bootstrap intervals resample users and
are descriptive conditional on fixed models, settings, and negative cohorts;
five r6.2 insiders provide particularly limited uncertainty information.

\paragraph{Changing scored targets.}
For a target-token class $c$, let $L_c(x)$ be its summed loss and $N_c(x)$ its
number of scored targets. Then
\begin{equation}
 s(x)=\frac{\sum_c L_c(x)}{\sum_c N_c(x)},\qquad
 s_B(x)=\frac{\sum_{c\in B}L_c(x)}{\sum_{c\in B}N_c(x)}.
\label{eq:behavior-score}
\end{equation}
Our behavior view includes \texttt{SESCOUNT}, \texttt{SES}, and any residual
non-profile target classes; the profile view includes \texttt{DAY} and
\texttt{PSY}. In this score-only comparison the full original input remains
visible to the model. A behavior target can therefore still be predicted from
profile information. Removing profile targets from an average, removing profile
\emph{input}, and retraining without profile content are three different
interventions.

\section{Opening the model}
\label{sec:method}
% claims: C02 C11 C12 C13 C14 C17 C19 C35 C38 C43
The interpretability method remains part of the work. Its purpose is to connect
an internal change to a measured output effect, while checking what the edit
actually changed. Its historical empirical claims require the repairs listed
later; the procedure itself does not depend on a favorable result.

\subsection{Adapt the model, then hold it fixed}
We adapt Qwen3-8B on non-insider user-days using QLoRA
\citep{qwen3,hu2022lora,dettmers2023qlora}. The base weights are frozen and
quantized; small low-rank matrices supply the learned updates. For a projection
matrix, the adapted weight is $W'=W+(\alpha/r)BA$. This makes it practical to
compare the same model with adapters enabled and disabled. The model is frozen
while we audit it. QLoRA supplies an accessible adaptation difference; it does
not establish that an LM is the best detector for these data.

\subsection{Isolate the adaptation difference}
At token position $t$ and layer $\ell$, compare the adapted and base states:
\begin{equation}
 \delta_t=h^{(\ell),\mathrm{ada}}_t-h^{(\ell),\mathrm{base}}_t.
\end{equation}
This is the representation difference induced by adaptation on that input.
It is not automatically a unique causal decomposition of the network.

\subsection{Find a sparse approximation}
Standardize each coordinate of $\delta_t$ and fit a sparse autoencoder (SAE):
\begin{equation}
 q_t=(\delta_t-\mu)\oslash\sigma,\qquad
 z_t=\operatorname{TopK}(\operatorname{ReLU}(Eq_t+b_e)),\qquad
 \widehat q_t=Dz_t+b_d.
\end{equation}
Only a few coefficients of $z_t$ are retained. The dictionary columns of $D$
are candidate directions of change, learned to reconstruct the observed deltas
\citep{cunningham2023sparse,gao2024scaling}. Sparse approximation is not a proof
that these directions are unique concepts. Any cross-dictionary comparison must
account for normalization: physical directions are $\sigma\odot D_{:,f}$.

\subsection{Select candidates, edit, and measure}
An activation difference between incident-window and reference examples can
rank candidate features. Selection is a hypothesis-generating step: labels,
reference-user eligibility, and the dictionary sample all affect the answer.
To test a selected set, edit its coefficients, continue the forward pass, and
measure the change in a declared output such as behavior-target NLL. Compare
against active features and random directions with the same edited support and
matched realized perturbation size. Keep a same-configuration unedited forward
pass as the reference. We measure both the requested edit and the change that
actually survives low-precision rounding.

A change in output establishes sensitivity to that specified edit. It does not
by itself prove that the features encode maliciousness, that the edit improves
anomaly ranking, or that the selected set is strictly necessary. Best-of-search
edits also differ from an intervention policy chosen before seeing its outcome.

\subsection{Locate activation without overinterpreting it}
For a fixed set of examples, activation mass at input class $c$ is
\begin{equation}
 m_f(c)=\frac{\sum_{x,t:\,\mathrm{class}(t)=c}z_{t,f}(x)}
                 {\sum_{x,t}z_{t,f}(x)}.
\end{equation}
This describes where a feature activates. A session-position hidden state can
use an earlier profile. High profile-position mass in selected features is
neither a complete account of the detector nor proof of a profile-only causal
mechanism. Corrected attribution must use the corrected selected features and
example population. The earlier profile-versus-behavior headline is pending
that recovery.

\section{Results supported after label correction}
\label{sec:findings}
\subsection{The comparison population changes the apparent detector quality}
% claims: C05 C06 C07 C08
\begin{center}\input{\evidenceroot/tables/recovery_detector.tex}\end{center}
The main-adapter day ROC changes from \rv{r42.all_users.day_roc} to
\rv{r42.unseen_users.day_roc} on r4.2 and from \rv{r62.all_users.day_roc} to
\rv{r62.unseen_users.day_roc} on r6.2 when negatives are restricted to unseen
users. This is a population/exposure comparison, not a controlled intervention
that identifies the sole cause of the gap. In r4.2, nine of the ten insiders
newly included by the label repair are scenario-3 users, and their full scores
rank incident-window days poorly. Their inclusion materially changes the main
result; recovery cannot stop at the previously covered 60 users.

The historical Deep SVDD, recurrent-autoencoder, and Isolation Forest rows are
not included as corrected comparisons. Their fitting populations need repair.
The corrected LM measurements do not establish a detector leaderboard ordering.

\subsection{Behavior-target scoring improves the complete r4.2 evaluation}
% claims: C32
\begin{center}\input{\evidenceroot/tables/recovery_behavior.tex}\end{center}
This separate, batch-size-one score decomposition covers all 42,468 unseen-pool
days, all 1,362 positive days, and all 70 insiders. A 256-example bridge exactly
reproduced the existing class-score pass before the missing users were scored.
Class losses reconstruct the total score, and an independent check reproduced
840 reported fold metrics to floating-point precision.

With the model and input held fixed, changing the average from all targets to
behavior targets raises mean fold day ROC from \rv{c32.full.day_roc} to
\rv{c32.behavior_only.day_roc}. The paired gain is \rv{c32.gain}, with a
descriptive 95\% user-bootstrap interval \rv{c32.gain.ci}. AP changes from
\rv{c32.full.day_ap} to \rv{c32.behavior_only.day_ap}. This supports a
\emph{score mitigation} for this model and population. It does not show that
profile information has stopped influencing behavior predictions.

Mean fold prevalence is \rv{c32.fold.prevalence}; pooled prevalence is
\rv{c32.pooled.prevalence}. The fold AP belongs with the former evaluation
construction, not the pooled prevalence. Neither AP value is the percentage
of alerts that would be correct at an operational threshold.
\begin{figure}[ht]
\centering\includegraphics[width=.94\textwidth]{\evidenceroot/figures/recovery_score_mitigation.pdf}
\caption{Complete r4.2 score decomposition. All three views retain the same input
and frozen model. The figure reports fold means, not pooled metrics.}
\label{fig:mitigation}
\end{figure}

\subsection{Training changes and score changes must stay separate}
% claims: C03 C09 C28 C31
The next table concerns separate r6.2 factorial adapters, not the main adapters
in the first table. Each input condition changes both adaptation and evaluation
content, with one adapter seed per condition.
\begin{center}\input{\evidenceroot/tables/recovery_factorial.tex}\end{center}
Removing all profile content improves ranking in both audited model families.
The 3B change is substantial in point estimate; the earlier description of a
uniformly tiny 3B effect is not supported. These are run-specific comparisons
on five insiders, not a general model-capacity law. Their training-byte
provenance is weaker than the main-adapter audit and remains explicitly qualified.

Scoring behavior targets while retaining the full input yields another comparison:
\begin{center}\input{\evidenceroot/tables/recovery_r62_behavior.tex}\end{center}
The 3B target-mask adapter omits profile targets from its training objective
while retaining profile context. Relative to behavior scoring of its full-training
counterpart, its day-ROC gain is about 0.006, with a descriptive interval crossing
zero, while AP decreases. This does not establish an additional mitigation gain.

Input masking of frozen main adapters is yet another experiment. Removing
personality input improves the cached day-ROC point estimates, but removing
all profile input does not have the same effect. Those arms used different
batch/hardware/cross-entropy configurations; the large personality-removal
contrast exceeds the observed metric sensitivity, while small differences
cannot be read at fine numerical precision. These findings are compatible with
several mechanisms and do not identify one by themselves.

\subsection{Temporal models provide a bounded comparison}
% claims: C33 C34
The corrected cached numeric-model comparison uses a common population of
141,662 history-eligible unseen-user days:
\begin{center}\input{\evidenceroot/tables/recovery_timeseries.tex}\end{center}
These are \emph{pooled} metrics. The seven-observation reconstruction includes
the target day; the forecast uses preceding observations only. Histories are
previous available observations, not necessarily consecutive calendar days.
All three checkpoint reloads passed the declared numerical tolerances, although
the seven-observation profile score used 99.5\% of its tolerance. That is a
tolerance pass, not bitwise equality. History-shuffle and broader copying or
mechanistic claims remain pending.

\section{What the causal evidence currently establishes}
\label{sec:causal-status}
% claims: C10 C11 C12 C13 C14 C15 C16 C17 C18 C20 C21 C22 C23 C24 C25 C26 C29 C35 C36 C37 C39
\subsection{The corrected pilot is an implementation diagnostic}
A corrected r4.2 pilot refit the sparse dictionary, reranked features using the
corrected discovery labels, and scored 158 incident-window/window-negative
pairs from 34 users under 15 conditions. The dictionary uses layer 26, 8,192
latents, and top-$k=4$; its sample, fit, and feature selection are recorded.
Live batch-size-one reference states and losses were recomputed. All 4,740
condition rows completed and the zero-edit controls reproduced the references.

One dictionary-direction control failed the predeclared 5\% realized-norm
matching tolerance at one window-negative token (6.42\% mismatch). Consequently
the original global gate remains failed. The following results are explicitly
\emph{post-run diagnostics}, with the original failed gate and all rows retained:
\begin{center}\input{\evidenceroot/tables/recovery_diagnostic.tex}\end{center}
The first two contrasts show larger behavior-loss increases for selected edits
than for their specified controls on incident-window receivers, but the
specificity contrast does not establish a
larger attack-specific response. These estimates do not reinstate the earlier
causal headline. Endpoint-local checks do not retroactively pass the global
protocol. A separate control-policy feasibility probe fixes the failing case
under a new sampling rule; it is not a replacement full experiment or a theorem
of scientific novelty.

\subsection{Saved responses retain a narrower meaning}
The historical Pilot 2 profile-swap measurements were relabeled on CPU. Of its
414 receivers, 129 change class; 86 have an ineligible partner in at least one
scored swap condition. The adapted model remains more sensitive than the base
model to the specified profile swaps in the frozen cohort. This measures input
sensitivity. Corrected class contrasts and eligible-partner subsets are post-hoc
descriptions; they do not restore the original matching or identify harm,
mediation, or a complete circuit.

The H6 direct-intervention pilot retains 1,320 scored rows for 120 receivers.
Forty-eight receiver labels change, and 44 receivers have at least one official
insider among their nominally benign profile partners. Its previously reported
small pooled mean effects reproduce. Small means for two specified lifts do not
establish equivalence of all feasible implementations or all users. The
corrected class summaries are likewise limited to the historical inputs and
selected dictionary.

The earlier input-swap experiment retains 723 eligible fixed r4.2 pairs and 60
r6.2 pairs after receiver and donor checks. There are 230 and 13 corrected
positive days absent from those original receiver caches, respectively. In
r4.2, some other cached receivers also need replacement donors. An extension
covering missing positives alone would therefore be incomplete.

\subsection{Historical claims awaiting dependency repair}
The old sparse-feature attribution percentages, repair/ablation estimates,
cross-dataset transfer, seed robustness, benign-only dictionary replication,
and smaller-model mechanistic replication are not presented as corrected
findings. Their dictionaries, ranking, negative pools, receivers, or donors can
depend on the faulty labels. For the repaired-3B study in particular, correcting
missing positive examples is insufficient: its label-selected negative sample
must also be redrawn. Historical artifacts are preserved for exact reuse checks,
not discarded; reuse requires the same complete intervention definition.

\section{Interpretation and remaining work}
\label{sec:discussion}
% claims: C19 C27 C30 C38 C40 C41 C42 C43 C44 C45 C46
The strongest corrected result is presently about the relationship between
\emph{evaluation population, score construction, and input dependence}. A model
can achieve a high ROC under one comparison population, perform much less well
on unseen users, and improve when only behavior targets contribute to its
score. Those observations motivate mechanistic investigation; they do not
already settle it.

The interpretability method remains useful: compare adaptation-induced states,
fit a sparse approximation, select candidates on a declared discovery group,
measure behavior-loss interventions against realized perturbation controls,
and report attribution with its location-versus-meaning limitation. Its
methodological value depends on executing these checks and exposing negative
or inconclusive outcomes. Existing sparse coding, causal contrasts, and
feasibility mathematics do not become new mathematics merely through application
to insider threats. The mathematical branch is retained as conditional work,
with no new theorem or novelty claim asserted here.

\paragraph{Comparison-model dependencies.}
The separate classical architecture probe fitted with ten r4.2 official
insiders and one r6.2 insider in its nominally normal training pool. For the
local session autoencoders, an audit of actual saved bundles found contaminated
training membership in 94/120 r4.2 bundles and 8/8 r6.2 bundles. All 120 r4.2
bundles have additional official insiders in their test groups. The 26 r4.2
bundles with clean training membership remain candidates for limited reuse;
they do not clear the old benchmark as a whole. The older sequence-baseline
harness shares the faulty fold construction and saves aggregate metrics,
without the checkpoints and per-example predictions needed for score-only
recovery. Refits must save scaler membership, checkpoints, predictions, and
corrected fold definitions before architecture comparisons return.

\paragraph{Scope outside CERT.}
The CERT clock bug does not resolve separate TWOS or LANL issues. The early
TWOS feature ranking had an independent defect; the corrected bounded held-out
study remains inconclusive. LANL has an independent coupling between sampling
and fold hashes. Neither is treated as a cleared replication here. Personality
novelty associations on only five r6.2 insiders remain exploratory and are not
used to explain the detector as a general law.

\paragraph{Next evidence needed.}
The next experiments must name the exact unsupported claim, corrected inputs,
frozen endpoints, eligibility checks, output artifacts, and measured runtime
before scaling. Priority goes to corrected causal selection and controls only
where the retained paper claim needs them; comparator refits must first repair
their training folds. Repeating every historical job is unnecessary, while
relabeling every old output is insufficient. The 46-family recovery ledger
records this distinction. This edition reports the verified subset and the
remaining gaps; it is not a claim that the complete mechanistic program is
finished or ready for submission.

\section{Conclusion}
Corrected labels preserve the main LM adapters but change several evaluated
populations and invalidate some label-dependent experiment choices. Complete
r4.2 behavior-target scoring improves ranking without changing the model or
removing profile context. The interpretability question therefore remains:
which internal pathways contribute to behavior predictions, and are their
measured effects specific to incident-window activity? The current diagnostic
pilot does not establish that specificity. A defensible account keeps this open
question visible while retaining the verified score-level findings.

\appendix
% Shared technical content keeps this edition complete with the main manuscript.
\section{Technical specification of the mechanistic audit}
\label{sec:technical}
% claims: C02 C11 C12 C13 C14 C15 C17 C19 C24 C35 C38 C43
\subsection{Adaptation and representation}
The main adapter configuration uses Qwen3-8B, NF4 four-bit base weights with
double quantization, bfloat16 computation, and rank-16 LoRA updates with scaling
32 and dropout 0.05 on attention and MLP projections. The declared adaptation
configuration uses maximum length 2,048, effective batch size eight, one epoch,
learning rate $1.5\times10^{-4}$, cosine decay, and 3\% warmup. Base and adapted
passes use the same input and tokenizer. Training configurations and artifact
hashes belong to each run; factorial adapters cannot be substituted for the
main adapters merely because their architecture is the same.

Let $\delta_t=h^{\mathrm{ada}}_t-h^{\mathrm{base}}_t$ at audited layer $\ell$.
With fitted location and scale vectors $\mu,\sigma$, the SAE minimizes
\begin{equation}
 \frac{1}{|\mathcal T|}\sum_{t\in\mathcal T}
 \left\|(\delta_t-\mu)\oslash\sigma-
       \bigl(D\operatorname{TopK}(\operatorname{ReLU}(E q_t+b_e))+b_d\bigr)\right\|_2^2,
 \qquad q_t=(\delta_t-\mu)\oslash\sigma.
\end{equation}
Both token sampling and optimization randomness must be recorded. The historical
SAE seed argument did not seed torch initialization; earlier across-seed changes
therefore mix sampling and optimization variability. The corrected fit records
both. A dictionary fitted on evaluation examples is transductive even if feature
ranking uses separate discovery users. It must not be described as a dictionary
learned independently of every evaluation user.

\subsection{Selection, attribution, and eligibility}
A mean activation-gap screen proposes a feature set,
\begin{equation}
 g_f=\mathbb E[z_f\mid y=1,\mathrm{discovery}]
       -\mathbb E[z_f\mid\mathrm{declared\ reference}],
 \qquad S=\operatorname{TopFive}(g).
\end{equation}
The reference population must be explicit: a window-negative insider day is not
interchangeable with a non-insider day. Active controls avoid inert features but
do not automatically match perturbation size. Token-weighted and equal-user
summaries also answer different questions. Attribution mass describes the stated
selected dictionary, token positions, and example set; changing any of these
requires recomputation. Physical dictionary directions are
$v_f=\operatorname{diag}(\sigma)D_{:,f}$, not standardized decoder columns alone.

\subsection{Two distinct intervention implementations}
Write $\widehat\delta_t=\mu+\sigma\odot(Dz_t+b_d)$ and choose a modified code
$z'_t$. A residual-preserving additive edit is
\begin{equation}
 h'_t=h^{\mathrm{ada}}_t+\sigma\odot D(z'_t-z_t).
\label{eq:additive}
\end{equation}
A reconstruction-replacement edit instead sets
\begin{equation}
 h'_t=h^{\mathrm{base}}_t+\mu+\sigma\odot(Dz'_t+b_d).
\label{eq:replacement}
\end{equation}
Their difference includes the original SAE reconstruction residual. Even if two
arms share replacement, downstream nonlinear responses need not cancel that
residual's effects. Unchanged coefficients, unedited tokens, and the code bias
must be handled consistently with the chosen implementation. Joint edits write
the declared coefficients over a declared union of token positions; their
interpretation is not a sum of proven single-feature effects.

For an output functional $s_V$ on declared targets $V$, define
$e_S(x)=s_V(x;I_S)-s_V(x;I_0)$, where $I_0$ is a live, same-configuration
reference. A selected-versus-control contrast is
$\mathbb E_u[\overline e_S(u)-\overline e_C(u)]$. An attack-specificity contrast
compares this within-user difference across incident-window and matched
window-negative receivers. Positive or negative signs must be interpreted using
this output definition; an increased NLL is not automatically better anomaly
ranking or successful mitigation.

Historical searched-donor repair used
\begin{equation}
 \tau=(\overline\Delta^S_{\rm anomalous}-\overline\Delta^S_{\rm noninsider})
       -(\overline\Delta^C_{\rm anomalous}-\overline\Delta^C_{\rm noninsider}),
\end{equation}
with the best signed score change retained over searched donors and strengths.
That defines a best-of-search estimand. A typical donor effect, a frozen
intervention policy, and detector AP after intervention require separate
measurements. The historical estimates are not reinstated in this edition.

\subsection{Realized perturbations and numerical consistency}
Let $r_t$ be the requested float32 update and $h_t$ the live low-precision state.
The realized perturbation is
\begin{equation}
 \widetilde r_t=operatorname{float32}
     (\operatorname{cast}(h_t+r_t))-operatorname{float32}(h_t).
\end{equation}
Matching $\|r_t\|$ does not guarantee matching $\|\widetilde r_t\|$. The corrected
pilot checks edited support and realized norm per token at a fixed 5\% relative
tolerance, alongside no-op identity and loss-partition checks. Norm matching is
one nuisance control; it does not make random directions semantically equivalent
to selected features.

The historical scoring passes differed with batch shape and cross-entropy
precision. Their per-example NLL variation, about 0.01 nats in the controlled
checks, can exceed the intervention contrasts of interest. The new reference
uses batch size one and float32 cross entropy. Cached hidden states are not
silently assumed equal to freshly computed states. Before reusing an effect,
check adapter, tokenizer/text, dictionary and normalization, feature set,
receiver, donor, token support, edit rule, score targets, and numerical
configuration. A change in any of these can invalidate exact reuse.

\subsection{Frozen cohorts and conditional mathematics}
Fixed-input reanalysis retains the originally measured responses and changes
only group labels. Post-hoc eligible-partner subsets have a different population
from the original sample. Within-user contrasts use users with both corrected
classes and do not recover the historical day pairing. All such estimates are
exploratory. No coefficient-feasibility statement alone proves output mediation:
a code-space target, its physically realizable hidden-state lift, and the model's
behavior under that lift are distinct objects. The mathematical branch remains
available for later work after these empirical dependencies are repaired.

\section{Recovery provenance and scope}
\label{app:provenance}
% claims: C01 C03 C04 C10 C19 C22 C27 C30 C35 C36 C38 C39 C40 C41 C42 C43 C44 C45 C46
The machine-readable source is \texttt{paper/data/recovery\_registry.json}.
It records all 46 claim families, source paths and SHA-256 digests, exact metric
selectors, completed scopes, limitations, and remaining work. Shared tables,
values, and the mitigation figure are generated from that registry for all
three editions. Source audits and the old manuscript texts are preserved in the
repair branch history and in the external recovery archive. Historical figures
and tables are not part of the corrected build.

\paragraph{Label validation.}
The official answer archive checksum was checked. Separate code enumerated the
master-window calendar target and compared all metadata rows; it did not import
the production label builder. Every named malicious event falls inside the
corresponding user's extracted session interval under the recorded calendar
clock. The r6.2 master includes a malformed CDE1846 start field, repaired using
that named user's earliest recorded event. Master/detail discrepancies for
ACM2278, CMP2946, and PLJ1771 are recorded rather than resolved by mixing target
definitions. Corrected evaluation JSONLs preserve all non-label values.

\paragraph{Score recovery.}
The cached registry preserves corrected fold means, pooled metrics, paired
contrasts, and fixed-sample results separately. Duplicate score identities must
agree before deduplication. Each score view must have finite values and valid
target counts; class losses must reconstruct total loss. The complete r4.2
extension follows an exact 256-row numerical bridge. The same-configuration
class pass and the older full-population pass differ slightly numerically, which
explains why the r4.2 full-score value in the mitigation table is 0.540 rather
than the first table's 0.538. They must not be spliced into one causal comparison.

\paragraph{New CPU recovery.}
Pilot 2 and H6 retain their saved losses and historical feature selection. An
independent arithmetic implementation reproduced 1,680 summaries and interval
endpoints within $6\times10^{-17}$; pooled H6 values also reproduce the earlier
independent audit. This checks aggregation, not the validity of the historical
sampling protocol. Donor manifests were checked against the official insider
list, independently of observed positive-day labels. The autoencoder audit reads
actual saved fold membership and hashes each bundle, rather than inferring
contamination from the manuscript's description.

\paragraph{Corrected causal pilot.}
The new r4.2 dictionary was fitted for 20 epochs on two million sampled token
deltas, with layer width 4,096 and 8,192 sparse features, top-$k=4$. Feature
ranking used corrected discovery labels. The dictionary remains transductive;
previously inspected confirmation users are not fresh holdouts. The full
158-pair run took approximately 45.9 minutes on an RTX 3070. Its original global
norm gate failed at one control token. Separate diagnostic analyses preserve
that failure, and the new conditional control-sampling feasibility probe is
recorded as a distinct policy and artifact. The runtime measures this pilot,
not the cost of all historical claim recovery.

\paragraph{Pending work and reusable artifacts.}
Raw sessions, corrected texts, main-adapter weights, and fixed-input responses
remain useful. A contaminated scaler or detector requires refitting; a changed
SAE dictionary requires downstream re-encoding, feature ranking, attribution,
and required interventions. The repaired-3B negative sample also needs correction.
The optional five-user personality-novelty association is deferred rather than
promoted as an explanation. TWOS, LANL, and conditional geometric results retain
their separate ledgers. The companion claim map and rerun specifications explain
what each retained or deferred result would require.

\section{Relation to earlier methods}
\label{app:related}
Anomaly detection supplies the one-class learning problem
\citep{chandola2009anomaly}; sequence and log models provide prior detector
approaches \citep{du2017deeplog,tuor2018recurrent,guo2021logbert}. Shortcut
learning and security-evaluation work motivate auditing the information behind
apparent predictive success \citep{geirhos2020shortcut,lapuschkin2019unmasking,arp2022dos,landauer2024critical}.
Sparse autoencoders provide candidate decompositions
\citep{cunningham2023sparse,gao2024scaling}, while activation-patching methodology
motivates matched interventions and careful interpretation
\citep{zhang2023best,heimersheim2024use}. These established methods are the
foundation of the audit. Applying them here does not itself establish a new
sparse-coding algorithm, causal-identification theorem, or superior detector.


\bibliography{../references}
\bibliographystyle{tmlr}
\end{document}
